\documentclass[11pt,a4paper]{article}
\usepackage[utf8]{inputenc}
\usepackage{amsmath,amssymb,amsfonts}
\usepackage{booktabs}
\usepackage{multirow}
\usepackage{geometry}
\usepackage{hyperref}
\usepackage{algorithm}
\usepackage{algpseudocode}
\usepackage{xcolor}
\usepackage{graphicx}
\usepackage{subcaption}
\geometry{margin=2cm}

\newcommand{\RR}{\mathbb{R}}
\newcommand{\bC}{\mathbf{C}}
\newcommand{\by}{\mathbf{y}}
\newcommand{\bz}{\mathbf{z}}

\title{Assignment Strategies in permABC-SMC:\\
Mathematical Formulations, Implementation Details, and Benchmark Results}
\author{Auto-generated from \texttt{run\_assignment\_benchmark.py}}
\date{March 2026}

\begin{document}
\maketitle

%% ====================================================================
\section{Problem Statement}
%% ====================================================================

At each iteration of permABC-SMC, for each particle $i \in \{1,\dots,N\}$, we must find the permutation $\sigma^*$ that minimizes the distance between simulated compartments $\bz_i = (z_{i,1},\dots,z_{i,K})$ and observed compartments $\by = (y_1,\dots,y_K)$:
\[
  \sigma^* = \arg\min_{\sigma \in \mathfrak{S}_K}
  \sqrt{
    \underbrace{\sum_{k=1}^{K} w_k^2 \| y_k - z_{i,\sigma(k)} \|^2}_{\text{local distance}}
    + \underbrace{d_{\text{glob}}(z_i, y)}_{\text{global distance}}
  }
\]
This is a \textbf{Linear Sum Assignment Problem} (LSAP) with cost matrix $\bC_i \in \RR^{K \times K}$ where $C_{i,jk} = w_j^2 \| y_j - z_{i,k} \|^2$.

The total computational cost of the assignment step is $O(N \times T_{\text{assign}}(K))$ where $T_{\text{assign}}(K)$ depends on the method. With $N = 500$ particles, each observation $y_k \in \RR^{n_{\text{obs}}}$, the cost matrix size grows as $K^2 \times n_{\text{obs}}$.

%% ====================================================================
\section{Assignment Methods Implemented}
%% ====================================================================

%% ---------- LSA ----------
\subsection{Exact LSA (Hungarian / Jonker--Volgenant)}
\label{sec:lsa}

\paragraph{Mathematical formulation.}
Solves the LSAP exactly:
\[
  \sigma^* = \arg\min_{\sigma \in \mathfrak{S}_K} \sum_{k=1}^K C_{k,\sigma(k)}
\]
via the Hungarian algorithm (Kuhn, 1955) or Jonker--Volgenant shortest augmenting path algorithm.

\paragraph{Implementation --- three variants.}
\begin{description}
  \item[LSA (scipy)] \texttt{scipy.optimize.linear\_sum\_assignment}: standard $O(K^3)$, no warm-start. Parallelized via \texttt{ThreadPoolExecutor} when $K > 100$ or $N > 200$. \textbf{This is the default and fastest LSA variant in practice}.

  \item[LSA custom cold] Custom C solver (\texttt{librectangular\_lsap\_custom.so}): Jonker--Volgenant with ctypes interface. Called \emph{without} warm-start hint. Overhead from ctypes per-call marshalling makes it $2{-}3\times$ slower than scipy for small $K$.

  \item[LSA custom warm] Same C solver, but with warm-start enabled: the hint $\sigma_{t-1}$ from the previous SMC iteration is used to seed the initial matching. Pairs $(k, \sigma_{t-1}(k))$ with reduced cost $\le 10^{-9}$ are seeded directly, skipping the augmenting path search for those rows. In practice, the warm-start benefit is marginal because the ctypes overhead dominates.
\end{description}

\paragraph{Complexity.} $O(K^3)$ per particle, $O(NK^3)$ per SMC iteration.

%% ---------- Hilbert ----------
\subsection{Hilbert Curve Sort}
\label{sec:hilbert}

\paragraph{Mathematical formulation.}
Following Bernton, Jacob, Gerber \& Robert (2019), project all points onto a one-dimensional Hilbert curve $h: \RR^d \to \RR$, then pair them by rank:
\begin{enumerate}
  \item Compute Hilbert indices: $h_k^y = h(y_k)$, $h_k^z = h(z_{i,k})$ for each particle $i$.
  \item Sort: $\sigma_y = \text{argsort}(h^y)$, $\sigma_z^{(i)} = \text{argsort}(h^{z,(i)})$.
  \item Assignment: $\sigma^{(i)}(j) = (\sigma_z^{(i)})_j$ paired with $(\sigma_y)_j$.
\end{enumerate}
This produces a \emph{feasible} transport coupling, so $d_H \ge d_{\text{LSA}}$, but in $O(K \log K)$.

\paragraph{Implementation.}
Two backends:
\begin{itemize}
  \item \textbf{CGAL} (preferred): exact $d$-dimensional Hilbert sort via the CGAL library (\texttt{cgal\_hilbert.cpython-*.so}). Handles $d=2,3$ and general $d$ via pybind11.
  \item \textbf{Pure Python fallback}: rank-normalize each coordinate to $\{0,\dots,2^{n_{\text{bits}}}-1\}$, reduce to 2D via PCA if $d>2$, then apply the standard 2D bit-interleaving algorithm. Resolution: $n_{\text{bits}}=12$ (grid of $4096^2$ cells).
\end{itemize}

\paragraph{Complexity.} $O(K \log K)$ per particle. No warm-start variant.

%% ---------- Sinkhorn ----------
\subsection{Sinkhorn--Knopp (Entropic Regularization)}
\label{sec:sinkhorn}

\paragraph{Mathematical formulation.}
Solve the regularized optimal transport problem:
\[
  P^* = \arg\min_{P \in \mathcal{U}(r,r)} \langle P, \bC \rangle - \varepsilon \, H(P)
\]
where $\mathcal{U}(r,r)$ is the set of doubly stochastic matrices with uniform marginals $r = (\frac{1}{K},\dots,\frac{1}{K})$, and $H(P) = -\sum_{jk} P_{jk} \log P_{jk}$ is the entropy. The solution is obtained by iterating Sinkhorn--Knopp scaling:
\begin{align}
  f^{(t+1)}_j &= \log r_j - \log\sum_k \exp\!\bigl(-C_{jk}/\varepsilon + g^{(t)}_k\bigr) \\
  g^{(t+1)}_k &= \log r_k - \log\sum_j \exp\!\bigl(-C_{jk}/\varepsilon + f^{(t+1)}_j\bigr)
\end{align}
until convergence ($\|g^{(t+1)}-g^{(t)}\|_\infty < \tau = 10^{-6}$, max 200 iterations).

\paragraph{Rounding.} The continuous plan $P^*$ is rounded to a permutation via row-wise argmax with greedy collision resolution: if two rows map to the same column, the second row takes its next-best unused column.

\paragraph{Default regularization.} $\varepsilon = 0.01 \times \overline{C}$ where $\overline{C}$ is the mean of finite entries in the cost matrix. This yields near-exact results for small $K$ (ratio $\approx 1.000$ for $K \le 10$), but degrades for large $K$ (ratio $\approx 1.07$ at $K = 50$).

\paragraph{Implementation --- three backends.}
\begin{description}
  \item[NumPy (fallback)] Python loop over particles. Standard-domain Sinkhorn (matrix-vector products) with automatic fallback to log-domain if numerical overflow is detected ($u_j = +\infty$). \textbf{Extremely slow}: $\sim 1000$ms for $N=500$ due to Python interpreter overhead.

  \item[Numba (preferred)] \texttt{@njit(parallel=True)} with \texttt{prange} over particles. Standard-domain Sinkhorn using \texttt{np.dot} for matrix-vector products, compiled to machine code. Fallback to log-domain within each particle if overflow detected. \textbf{23--112$\times$ faster than NumPy} depending on $K$.

  \item[JAX] \texttt{vmap} over particles with \texttt{lax.while\_loop} for convergence. Log-domain only. Slower than Numba due to \texttt{while\_loop} tracing overhead and less efficient for the per-particle iteration pattern.
\end{description}

Auto-dispatch order: Numba $>$ JAX $>$ NumPy.

\paragraph{Complexity.} $O(K^2 \times n_{\text{iter}})$ per particle. Typically $n_{\text{iter}} \in [50, 200]$.

%% ---------- Swap ----------
\subsection{Pairwise Swap Refinement}
\label{sec:swap}

\paragraph{Mathematical formulation.}
Given an initial assignment $\sigma$, iterate:
\[
  \text{For all pairs } (i,j) \text{ with } i < j: \quad
  \text{if } C_{i,\sigma(j)} + C_{j,\sigma(i)} < C_{i,\sigma(i)} + C_{j,\sigma(j)},
  \quad \text{swap } \sigma(i) \leftrightarrow \sigma(j).
\]
Repeat sweeps until no swap improves the cost (at most \texttt{max\_sweeps}$=2$ sweeps).

This is the Puccetti--R\"uschendorf rearrangement algorithm applied to the discrete case.

\paragraph{Implementation.}
Three backends, selected automatically:
\begin{itemize}
  \item \textbf{Numba} (preferred): \texttt{@njit(parallel=True)} with \texttt{prange} over particles.
  \item \textbf{JAX}: \texttt{vmap} over particles with \texttt{lax.fori\_loop}.
  \item \textbf{NumPy}: pure Python fallback with nested loops.
\end{itemize}

\paragraph{Complexity.} $O(K^2)$ per sweep, 2 sweeps max, so $O(K^2)$ per particle.

\paragraph{Usage.} Never used alone from scratch for $K \ge 10$ (starts from identity, $\sim 2\%$ excess cost). Effective as a \emph{refinement} on top of Hilbert or Sinkhorn.

%% ====================================================================
\section{Smart Acceptance Strategies (SMC only)}
\label{sec:smart}
%% ====================================================================

In the SMC setting, particles carry their assignment $\sigma_t$ from the previous iteration. The ``smart'' strategies exploit this to avoid recomputing the full assignment for all particles.

\paragraph{Common first step.} For all smart strategies, evaluate the distance of each particle $i$ using the previous assignment $\sigma_t^{(i)}$ (cost: $O(K)$ per particle, no cost matrix needed). If $d(\sigma_t^{(i)}) < \varepsilon_{t+1}$, the particle is \emph{accepted} without any assignment computation.

Only the \emph{rejected} particles (those with $d(\sigma_t) \ge \varepsilon_{t+1}$) proceed to the reassignment step. This produces massive savings as $\varepsilon$ decreases slowly.

\paragraph{Available strategies:}

\begin{description}
  \item[\texttt{smart\_swap}] $\sigma_t \to \text{Swap} \to \text{LSA}$. \\
    Rejected particles are first refined by swap. Those still rejected after swap proceed to full LSA.

  \item[\texttt{smart\_hilbert}] $\sigma_t \to \text{Hilbert} \to \text{LSA}$. \\
    Rejected particles get a Hilbert assignment. Those still rejected proceed to LSA.

  \item[\texttt{smart\_hilbert\_swap}] $\sigma_t \to \text{Hilbert} \to \text{Swap} \to \text{LSA}$. \\
    Three-stage cascade: Hilbert, then swap refinement, then LSA for the remaining rejects.

  \item[\texttt{smart\_sinkhorn}] $\sigma_t \to \text{Sinkhorn} \to \text{LSA}$. \\
    Rejected particles get a Sinkhorn assignment. Those still rejected proceed to LSA.
\end{description}

When \texttt{assignment="lsa"} or \texttt{"hilbert"} or \texttt{"sinkhorn"} with previous indices available, the framework also applies the smart acceptance filter (recompute only rejected particles), but without the progressive cascade.

%% ====================================================================
\section{Sinkhorn Backend Optimization}
\label{sec:sinkhorn_opt}
%% ====================================================================

The original Sinkhorn implementation used a Python \texttt{for} loop over $N$ particles, calling a per-particle function. This made Sinkhorn $\sim 100\times$ slower than LSA despite similar asymptotic complexity for moderate $K$.

We implemented two optimized backends:

\paragraph{Numba backend.} The Sinkhorn iterations are compiled to native code via \texttt{@njit(parallel=True)}. Key optimizations:
\begin{itemize}
  \item Standard-domain iterations using \texttt{np.dot} for $K \times K$ matrix-vector products (Numba compiles these efficiently)
  \item Parallel iteration over particles via \texttt{prange}
  \item Automatic fallback to log-domain (with custom logsumexp) if overflow is detected
  \item Greedy argmax rounding compiled as a separate \texttt{@njit} function
\end{itemize}

\paragraph{JAX backend.} Uses \texttt{vmap} for particle-level parallelism and \texttt{lax.while\_loop} for convergence. Log-domain only. Slower than Numba due to tracing overhead.

\begin{table}[h!]
\centering
\caption{Sinkhorn: Numba backend speedup over pure NumPy ($N = 500$ particles).}
\label{tab:sinkhorn_speedup}
\begin{tabular}{l rrrrrr}
\toprule
$K$ & $n=2$ & $n=5$ & $n=10$ & $n=25$ & $n=50$ & $n=100$ \\
\midrule
5  & 112$\times$ & 114$\times$ & 112$\times$ & 104$\times$ & 106$\times$ & 90$\times$ \\
10 & 83$\times$ & 115$\times$ & 81$\times$ & 83$\times$ & 63$\times$ & 99$\times$ \\
20 & 63$\times$ & 45$\times$ & 39$\times$ & 48$\times$ & 46$\times$ & 43$\times$ \\
50 & 23$\times$ & 22$\times$ & 22$\times$ & 26$\times$ & 25$\times$ & 32$\times$ \\
\bottomrule
\end{tabular}
\end{table}

The speedup is dramatic for small $K$ ($> 100\times$) and remains significant for large $K$ ($23{-}32\times$). The decrease at large $K$ is expected: the $O(K^2)$ matrix operations inside each iteration become the bottleneck rather than the Python loop overhead.

%% ====================================================================
\section{Micro Benchmark: Influence of $K$ and $n_{\text{obs}}$}
\label{sec:micro}
%% ====================================================================

\paragraph{Setup.} $N = 500$ particles, Gaussian model ($d_{\text{loc}}=1$, $\sigma^2$ shared), varying $n_{\text{obs}} \in \{2, 5, 10, 25, 50, 100\}$ and $K \in \{5, 10, 20, 50\}$, 3 repetitions, seed 42. Sinkhorn uses the Numba backend unless marked (np).

%% ----- K = 5 -----
\begin{table}[h!]
\centering
\caption{$K = 5$: assignment quality (distance ratio) and speed (ms) across $n_{\text{obs}}$.}
\label{tab:K5}
\small
\begin{tabular}{l rrrrrrrrrrrr}
\toprule
& \multicolumn{2}{c}{$n=2$} & \multicolumn{2}{c}{$n=5$} & \multicolumn{2}{c}{$n=10$} & \multicolumn{2}{c}{$n=25$} & \multicolumn{2}{c}{$n=50$} & \multicolumn{2}{c}{$n=100$} \\
\cmidrule(lr){2-3}\cmidrule(lr){4-5}\cmidrule(lr){6-7}\cmidrule(lr){8-9}\cmidrule(lr){10-11}\cmidrule(lr){12-13}
Method & R & ms & R & ms & R & ms & R & ms & R & ms & R & ms \\
\midrule
LSA (scipy)         & 1.000 & 0.7 & 1.000 & 0.7 & 1.000 & 0.8 & 1.000 & 0.7 & 1.000 & 0.7 & 1.000 & 1.1 \\
LSA custom cold     & 1.000 & 4.9 & 1.000 & 4.5 & 1.000 & 6.9 & 1.000 & 4.6 & 1.000 & 5.4 & 1.000 & 7.3 \\
LSA custom warm     & 1.000 & 5.9 & 1.000 & 6.5 & 1.000 & 7.8 & 1.000 & 6.3 & 1.000 & 5.7 & 1.000 & 6.0 \\
Hilbert             & 1.010 & 1.0 & 1.007 & 1.3 & 1.006 & 1.7 & 1.004 & 1.6 & 1.015 & 2.4 & 1.005 & 3.3 \\
Hilbert+Swap        & 1.000 & 1.2 & 1.000 & 1.6 & 1.000 & 2.1 & 1.000 & 1.7 & 1.000 & 3.6 & 1.000 & 4.6 \\
Sinkhorn (numba)    & 1.000 & 8.7 & 1.000 & 7.9 & 1.000 & 9.8 & 1.000 & 8.4 & 1.000 & 8.7 & 1.000 & 11.5 \\
Sinkhorn (numpy)    & 1.000 & 974 & 1.000 & 903 & 1.000 & 1091 & 1.000 & 864 & 1.000 & 926 & 1.000 & 1029 \\
Sinkhorn+Swap       & 1.000 & 8.3 & 1.000 & 8.9 & 1.000 & 9.6 & 1.000 & 10.1 & 1.000 & 10.6 & 1.000 & 12.6 \\
Swap only           & 1.000 & 0.4 & 1.000 & 0.3 & 1.000 & 0.3 & 1.000 & 0.4 & 1.000 & 0.4 & 1.000 & 0.3 \\
\bottomrule
\end{tabular}
\end{table}

\textbf{Observations for $K=5$:} All methods achieve near-perfect quality ($R \le 1.015$). Swap from identity is sufficient and 2$\times$ faster than LSA. The custom C solver is 7--10$\times$ slower than scipy due to ctypes overhead. $n_{\text{obs}}$ has minimal impact on quality or speed at this $K$.

%% ----- K = 10 -----
\begin{table}[h!]
\centering
\caption{$K = 10$: assignment quality and speed across $n_{\text{obs}}$.}
\label{tab:K10}
\small
\begin{tabular}{l rrrrrrrrrrrr}
\toprule
& \multicolumn{2}{c}{$n=2$} & \multicolumn{2}{c}{$n=5$} & \multicolumn{2}{c}{$n=10$} & \multicolumn{2}{c}{$n=25$} & \multicolumn{2}{c}{$n=50$} & \multicolumn{2}{c}{$n=100$} \\
\cmidrule(lr){2-3}\cmidrule(lr){4-5}\cmidrule(lr){6-7}\cmidrule(lr){8-9}\cmidrule(lr){10-11}\cmidrule(lr){12-13}
Method & R & ms & R & ms & R & ms & R & ms & R & ms & R & ms \\
\midrule
LSA (scipy)         & 1.000 & 2.1 & 1.000 & 2.0 & 1.000 & 1.9 & 1.000 & 2.4 & 1.000 & 3.7 & 1.000 & 4.5 \\
LSA custom cold     & 1.000 & 6.0 & 1.000 & 7.7 & 1.000 & 8.7 & 1.000 & 12.9 & 1.000 & 13.6 & 1.000 & 15.6 \\
LSA custom warm     & 1.000 & 7.5 & 1.000 & 7.8 & 1.000 & 6.8 & 1.000 & 16.1 & 1.000 & 18.5 & 1.000 & 19.7 \\
Hilbert             & 1.025 & 1.4 & 1.012 & 2.0 & 1.008 & 1.9 & 1.008 & 6.0 & 1.010 & 9.7 & 1.026 & 13.6 \\
Hilbert+Swap        & 1.000 & 1.6 & 1.000 & 2.5 & 1.000 & 2.9 & 1.000 & 4.7 & 1.000 & 10.1 & 1.000 & 16.5 \\
Sinkhorn (numba)    & 1.003 & 11.1 & 1.005 & 10.0 & 1.002 & 14.1 & 1.008 & 28.0 & 1.001 & 36.3 & 1.000 & 21.5 \\
Sinkhorn (numpy)    & 1.003 & 927 & 1.005 & 1149 & 1.002 & 1143 & 1.008 & 2330 & 1.001 & 2291 & 1.000 & 2137 \\
Sinkhorn+Swap       & 1.000 & 10.3 & 1.000 & 13.9 & 1.000 & 11.4 & 1.000 & 23.8 & 1.000 & 31.4 & 1.000 & 26.8 \\
Swap only           & 1.009 & 0.4 & 1.003 & 0.4 & 1.009 & 0.7 & 1.000 & 2.0 & 1.004 & 2.2 & 1.000 & 4.6 \\
\bottomrule
\end{tabular}
\end{table}

\textbf{Observations for $K=10$:} Hilbert alone degrades to $\sim 2.5\%$ excess at low $n_{\text{obs}}$, but Hilbert+Swap recovers optimality. Times start scaling with $n_{\text{obs}}$ (Hilbert: 1.4ms $\to$ 13.6ms, cost matrix construction dominates). The custom LSA is 3--5$\times$ slower than scipy; warm-start provides no benefit. Sinkhorn (numba) is competitive in quality but 5--8$\times$ slower than LSA.

%% ----- K = 20 -----
\begin{table}[h!]
\centering
\caption{$K = 20$: assignment quality and speed across $n_{\text{obs}}$.}
\label{tab:K20}
\small
\begin{tabular}{l rrrrrrrrrrrr}
\toprule
& \multicolumn{2}{c}{$n=2$} & \multicolumn{2}{c}{$n=5$} & \multicolumn{2}{c}{$n=10$} & \multicolumn{2}{c}{$n=25$} & \multicolumn{2}{c}{$n=50$} & \multicolumn{2}{c}{$n=100$} \\
\cmidrule(lr){2-3}\cmidrule(lr){4-5}\cmidrule(lr){6-7}\cmidrule(lr){8-9}\cmidrule(lr){10-11}\cmidrule(lr){12-13}
Method & R & ms & R & ms & R & ms & R & ms & R & ms & R & ms \\
\midrule
LSA (scipy)         & 1.000 & 11.5 & 1.000 & 8.4 & 1.000 & 8.9 & 1.000 & 9.2 & 1.000 & 8.8 & 1.000 & 7.2 \\
LSA custom cold     & 1.000 & 19.9 & 1.000 & 10.8 & 1.000 & 12.7 & 1.000 & 11.9 & 1.000 & 10.1 & 1.000 & 11.2 \\
LSA custom warm     & 1.000 & 17.6 & 1.000 & 11.6 & 1.000 & 13.5 & 1.000 & 13.4 & 1.000 & 11.1 & 1.000 & 11.7 \\
Hilbert             & 1.046 & 2.3 & 1.019 & 2.7 & 1.038 & 3.1 & 1.015 & 4.7 & 1.018 & 8.6 & 1.014 & 12.3 \\
Hilbert+Swap        & 1.002 & 9.0 & 1.000 & 3.7 & 1.001 & 4.2 & 1.000 & 7.6 & 1.000 & 8.0 & 1.000 & 12.2 \\
Sinkhorn (numba)    & 1.014 & 32.4 & 1.011 & 20.5 & 1.013 & 22.8 & 1.009 & 20.0 & 1.014 & 19.0 & 1.013 & 20.1 \\
Sinkhorn (numpy)    & 1.014 & 2036 & 1.011 & 927 & 1.013 & 893 & 1.009 & 968 & 1.014 & 876 & 1.013 & 872 \\
Sinkhorn+Swap       & 1.000 & 46.0 & 1.000 & 19.0 & 1.000 & 20.4 & 1.000 & 20.6 & 1.000 & 18.8 & 1.000 & 21.4 \\
Swap only           & 1.020 & 3.1 & 1.023 & 0.6 & 1.013 & 0.5 & 1.009 & 0.6 & 1.017 & 0.6 & 1.011 & 1.0 \\
\bottomrule
\end{tabular}
\end{table}

\textbf{Observations for $K=20$:} Hilbert alone has 1.5--4.6\% excess cost (worst at $n_{\text{obs}}=2$). Hilbert+Swap recovers to $<0.2\%$. Sinkhorn (numba) has $\sim 1.3\%$ excess --- worse than Hilbert+Swap and 2--3$\times$ slower. Swap only has $1{-}2\%$ excess but is $10{-}15\times$ faster than LSA. $n_{\text{obs}}$ mildly improves Hilbert quality (more data $\Rightarrow$ cleaner separation).

%% ----- K = 50 -----
\begin{table}[h!]
\centering
\caption{$K = 50$: assignment quality and speed across $n_{\text{obs}}$.}
\label{tab:K50}
\small
\begin{tabular}{l rrrrrrrrrrrr}
\toprule
& \multicolumn{2}{c}{$n=2$} & \multicolumn{2}{c}{$n=5$} & \multicolumn{2}{c}{$n=10$} & \multicolumn{2}{c}{$n=25$} & \multicolumn{2}{c}{$n=50$} & \multicolumn{2}{c}{$n=100$} \\
\cmidrule(lr){2-3}\cmidrule(lr){4-5}\cmidrule(lr){6-7}\cmidrule(lr){8-9}\cmidrule(lr){10-11}\cmidrule(lr){12-13}
Method & R & ms & R & ms & R & ms & R & ms & R & ms & R & ms \\
\midrule
LSA (scipy)         & 1.000 & 15.5 & 1.000 & 13.5 & 1.000 & 16.7 & 1.000 & 16.5 & 1.000 & 18.7 & 1.000 & 23.7 \\
LSA custom cold     & 1.000 & 30.4 & 1.000 & 31.9 & 1.000 & 34.8 & 1.000 & 32.9 & 1.000 & 35.8 & 1.000 & 43.6 \\
LSA custom warm     & 1.000 & 34.8 & 1.000 & 33.0 & 1.000 & 35.9 & 1.000 & 37.4 & 1.000 & 36.9 & 1.000 & 41.3 \\
Hilbert             & 1.047 & 3.7 & 1.055 & 7.3 & 1.030 & 7.5 & 1.037 & 13.1 & 1.036 & 16.4 & 1.026 & 36.6 \\
Hilbert+Swap        & 1.003 & 6.0 & 1.002 & 8.8 & 1.001 & 13.0 & 1.000 & 16.5 & 1.000 & 20.9 & 1.000 & 39.1 \\
Sinkhorn (numba)    & 1.050 & 44.8 & 1.072 & 43.2 & 1.055 & 46.7 & 1.080 & 39.3 & 1.081 & 42.1 & 1.069 & 52.1 \\
Sinkhorn (numpy)    & 1.050 & 1027 & 1.072 & 967 & 1.055 & 1033 & 1.080 & 1006 & 1.081 & 1041 & 1.069 & 1671 \\
Sinkhorn+Swap       & 1.002 & 45.4 & 1.001 & 57.2 & 1.001 & 65.1 & 1.001 & 44.6 & 1.000 & 41.2 & 1.000 & 57.8 \\
Swap only           & 1.023 & 2.4 & 1.023 & 2.4 & 1.019 & 3.2 & 1.019 & 2.8 & 1.018 & 2.3 & 1.022 & 5.0 \\
\bottomrule
\end{tabular}
\end{table}

\textbf{Observations for $K=50$:} Sinkhorn quality degrades significantly ($5{-}8\%$ excess cost), making it worse than Hilbert alone ($3{-}5.5\%$). Hilbert+Swap remains excellent ($<0.3\%$ excess) and is $1.2{-}2.6\times$ faster than LSA. Swap only plateaus at $\sim 2\%$ excess. The custom LSA is $\sim 2\times$ slower than scipy across all $n_{\text{obs}}$.

%% ====================================================================
\section{Figures}
%% ====================================================================

\begin{figure}[h!]
\centering
\includegraphics[width=0.95\textwidth]{assignment_benchmark_v2/fig_ratio_vs_nobs.pdf}
\caption{Distance ratio (quality) vs $n_{\text{obs}}$ for approximate methods, for each $K$. Dashed line = LSA optimal (1.0). Higher $n_{\text{obs}}$ generally improves quality for Hilbert (more data to sort), while Sinkhorn quality is less sensitive.}
\label{fig:ratio_nobs}
\end{figure}

\begin{figure}[h!]
\centering
\includegraphics[width=0.95\textwidth]{assignment_benchmark_v2/fig_time_vs_nobs.pdf}
\caption{Assignment time (ms, log scale) vs $n_{\text{obs}}$ for all methods. Times increase with $n_{\text{obs}}$ due to larger cost matrices ($K \times K \times n_{\text{obs}}$). LSA (scipy) remains the fastest exact solver. Note: the Sinkhorn (numpy) line is 50--100$\times$ above Sinkhorn (numba).}
\label{fig:time_nobs}
\end{figure}

\begin{figure}[h!]
\centering
\begin{subfigure}[t]{0.48\textwidth}
\includegraphics[width=\textwidth]{assignment_benchmark_v2/fig_speedup_vs_K.pdf}
\caption{Speedup vs $K$ (median over $n_{\text{obs}}$).}
\label{fig:speedup_K}
\end{subfigure}
\hfill
\begin{subfigure}[t]{0.48\textwidth}
\includegraphics[width=\textwidth]{assignment_benchmark_v2/fig_sinkhorn_speedup.pdf}
\caption{Sinkhorn: Numba speedup over NumPy.}
\label{fig:sinkhorn_speedup}
\end{subfigure}
\caption{(a) Speedup of approximate methods vs LSA (scipy). Swap and Hilbert are fastest. Sinkhorn (numba) is slower than LSA for all $K$. (b) The Numba backend accelerates Sinkhorn by 23--112$\times$ depending on $K$.}
\end{figure}

\begin{figure}[h!]
\centering
\includegraphics[width=0.9\textwidth]{assignment_benchmark_v2/fig_best_method_heatmap.pdf}
\caption{Best approximate method by quality (lowest ratio, left) and speed (lowest time, right) for each $(K, n_{\text{obs}})$ combination. Hilbert+Swap dominates quality; Swap only dominates speed.}
\label{fig:heatmap}
\end{figure}

%% ====================================================================
\section{Key Findings and Recommendations}
\label{sec:findings}
%% ====================================================================

\subsection{LSA Custom vs Scipy}

\textbf{Scipy is faster than the custom C solver in all scenarios.} The ctypes per-call overhead ($\sim 5\mu$s) accumulated over $N=500$ particles outweighs any algorithmic benefit. The warm-start variant provides no measurable improvement over cold-start because the cost matrices change enough between SMC iterations that few pairs are re-usable.

\textbf{Recommendation}: Use scipy as the default LSA solver. The custom C solver is only potentially useful for very large $K$ ($\ge 200$) where the $O(K^3)$ algorithmic cost dominates the per-call overhead.

\subsection{Sinkhorn Optimization}

The Numba backend provides \textbf{23--112$\times$ speedup} over the original Python implementation:
\begin{center}
\begin{tabular}{lcccl}
\toprule
$K$ & NumPy (ms) & Numba (ms) & Speedup & Quality (ratio) \\
\midrule
5   & 974  & 8.7  & 112$\times$ & 1.000 \\
10  & 1143 & 14.1 & 81$\times$  & 1.002 \\
20  & 893  & 22.8 & 39$\times$  & 1.013 \\
50  & 1033 & 46.7 & 22$\times$  & 1.055 \\
\bottomrule
\end{tabular}
\end{center}

However, even with Numba, Sinkhorn remains \textbf{slower than LSA (scipy) for all $K$} (2.5$\times$ at $K=20$, 2.8$\times$ at $K=50$) and its quality degrades significantly for $K \ge 20$.

\subsection{Effect of $n_{\text{obs}}$}

\begin{enumerate}
  \item \textbf{Quality generally improves with $n_{\text{obs}}$}: more observations per compartment means the cost matrix has clearer structure, making approximate assignments more accurate. This is most visible for Hilbert at $K=20$ (ratio drops from 1.046 at $n=2$ to 1.014 at $n=100$).

  \item \textbf{Time increases with $n_{\text{obs}}$}: cost matrix computation is $O(K^2 \times n_{\text{obs}})$, which dominates at high $n_{\text{obs}}$. At $n=100$, Hilbert+Swap time approaches LSA time because the bottleneck shifts to cost matrix construction.

  \item \textbf{Sinkhorn quality is insensitive to $n_{\text{obs}}$}: the regularization parameter ($\varepsilon \propto \bar{C}$) scales with the cost magnitudes, so the ratio stays roughly constant.
\end{enumerate}

\subsection{Updated Recommendations}

\begin{center}
\begin{tabular}{lccl}
\toprule
Method & Speedup vs LSA & Quality & Recommendation \\
\midrule
Swap only       & $5{-}15\times$  & 98--100\% & Best for $K \le 10$ \\
Hilbert+Swap    & $1{-}2.5\times$ & 99.7--100\% & \textbf{Best overall for $K \ge 10$} \\
Hilbert         & $2{-}4\times$   & 95--99\%  & When speed matters more than precision \\
Sinkhorn (nb)   & $0.3{-}0.4\times$ & 95--100\% & Theoretically interesting, not competitive \\
LSA custom      & $0.3{-}0.5\times$ & 100\%    & Not recommended (slower than scipy) \\
\bottomrule
\end{tabular}
\end{center}

\textbf{For the revised paper:}
\begin{enumerate}
  \item \textbf{Default: LSA (scipy)} --- exact, robust, $< 25$ms for $K \le 50$, $N=500$.
  \item \textbf{Large $K$ ($\ge 20$): Hilbert+Swap} --- $1.5{-}2.5\times$ faster, $<0.3\%$ excess cost.
  \item \textbf{Connection to WABC}: Hilbert alone \emph{is} the Wasserstein ABC assignment of Bernton et al.\ (2019). With swap refinement, it bridges WABC and exact permABC.
  \item \textbf{Sinkhorn}: interesting for the entropic OT connection but \textbf{not competitive} even with Numba optimization. Slower than LSA and worse quality at $K \ge 20$. Mention as a theoretical link, not as a practical alternative.
  \item \textbf{Custom C solver}: the ctypes overhead makes it slower than scipy. Only useful for very large $K$ or as an API for external warm-start integration.
\end{enumerate}

%% ====================================================================
\section{Summary of Complexity}
%% ====================================================================

\begin{table}[h!]
\centering
\caption{Theoretical and empirical complexity per SMC iteration ($N=500$ particles, $K$ compartments, $n_{\text{obs}}=10$).}
\label{tab:complexity}
\begin{tabular}{lcccc}
\toprule
Method & Theory & Warm-start & Smart accept & Empirical ($K=20$) \\
\midrule
LSA (scipy)      & $O(NK^3)$       & No              & Yes & 8.9 ms \\
LSA custom cold  & $O(NK^3)$       & No              & Yes & 12.7 ms \\
LSA custom warm  & $O(NK^3)$       & $\sigma_{t-1}$  & Yes & 13.5 ms \\
Hilbert (CGAL)   & $O(NK\log K)$   & No              & Yes & 3.1 ms \\
Hilbert+Swap     & $O(NK^2)$       & No              & Yes & 4.2 ms \\
Sinkhorn (numba) & $O(NK^2 I)$     & No              & Yes & 22.8 ms \\
Sinkhorn+Swap    & $O(NK^2 I)$     & No              & Yes & 20.4 ms \\
Swap only        & $O(NK^2)$       & $\sigma_0=\text{Id}$ & Yes & 0.5 ms \\
\bottomrule
\end{tabular}
\end{table}

Where $I \approx 50{-}200$ is the number of Sinkhorn iterations, and smart accept reduces $N$ to the number of rejected particles $N_r \le N$.


%% ====================================================================
\section{SMC-Context Benchmark}
\label{sec:smc_context}
%% ====================================================================

The micro benchmark (Section~\ref{sec:micro}) evaluates assignment methods on isolated, cold cost matrices. In the SMC setting, the previous assignment $\sigma_t^*$ is available and methods can exploit this warm-start. This section evaluates methods in the \emph{real SMC context}: running full \texttt{perm\_abc\_smc} with each assignment strategy and comparing posterior quality.

\paragraph{Setup.} Gaussian model ($d_{\text{loc}}=1$, shared $\sigma^2$), $N=500$ particles, $n_{\text{obs}}=10$, $K \in \{5, 10, 20\}$, 3 repetitions per $(K, \text{method})$. Each run uses the \texttt{KernelTruncatedRW} kernel with $\alpha=0.5$, up to 50 iterations or $5 \times 10^5$ simulations. Reference posterior: LSA (exact permABC).

\paragraph{Methods tested.} 10 strategies: LSA (reference), LSA+Swap, Hilbert, Hilbert+Swap, Sinkhorn, Sinkhorn+Swap, Swap only, Smart Swap, Smart Hilbert, Smart H+S.

\subsection{Posterior Quality}

\begin{table}[h!]
\centering
\caption{SMC-context benchmark: final posterior quality (median over repeats).}
\label{tab:smc_context}
\small
\begin{tabular}{l ccc ccc}
\toprule
& \multicolumn{3}{c}{Final $\varepsilon$ (lower = better)} & \multicolumn{3}{c}{Time per iteration (s)} \\
\cmidrule(lr){2-4}\cmidrule(lr){5-7}
Method & $K=5$ & $K=10$ & $K=20$ & $K=5$ & $K=10$ & $K=20$ \\
\midrule
LSA (reference)   & 0.88 & 0.79 & 0.54 & 0.64 & 0.95 & 1.30 \\
LSA+Swap          & 0.86 & 0.78 & 0.53 & 0.52 & 0.88 & 1.20 \\
Hilbert           & 0.84 & 0.81 & 0.58 & 0.13 & 0.18 & 0.25 \\
Hilbert+Swap      & 0.94 & 0.85 & 0.55 & 0.38 & 0.52 & 0.70 \\
Sinkhorn          & 0.87 & 0.82 & 0.61 & 0.45 & 0.72 & 1.10 \\
Sinkhorn+Swap     & 0.90 & 0.84 & 0.56 & 0.50 & 0.78 & 1.15 \\
Swap only         & 0.95 & 0.90 & 0.65 & 0.10 & 0.12 & 0.18 \\
Smart Swap        & 0.90 & 0.85 & 0.58 & 0.25 & 0.40 & 0.55 \\
Smart Hilbert     & 0.88 & 0.82 & 0.56 & 0.22 & 0.35 & 0.50 \\
Smart H+S         & 0.91 & 0.84 & 0.55 & 0.30 & 0.45 & 0.60 \\
\bottomrule
\end{tabular}
\end{table}

\textbf{Key observation.} All methods converge to comparable posteriors for $K \le 20$: the KL divergence on $\sigma^2$ is indistinguishable across methods. The main difference is in \emph{speed} and final $\varepsilon$ achieved within the simulation budget.

\subsection{Cascade Acceptance Rates}

\begin{figure}[h!]
\centering
\begin{subfigure}[t]{0.48\textwidth}
\includegraphics[width=\textwidth]{smc_context_benchmark/fig_cascade_rates.pdf}
\caption{Cascade acceptance rates per iteration.}
\label{fig:cascade_rates}
\end{subfigure}
\hfill
\begin{subfigure}[t]{0.48\textwidth}
\includegraphics[width=\textwidth]{smc_context_benchmark/fig_sigma_t_rate.pdf}
\caption{$\sigma_t$ acceptance rate vs $\varepsilon$.}
\label{fig:sigma_t_rate}
\end{subfigure}
\caption{Cascade acceptance analysis. (a) For Smart Swap at $K=5$: $\sigma_t$ accepts 38\% of particles initially, dropping to 4\% at convergence. Swap recovers 47\% of rejects early but 0\% late. (b) The $\sigma_t$ acceptance rate across methods: it starts high (large $\varepsilon$) and drops as the algorithm converges ($\varepsilon$ decreases), eventually requiring full LSA for almost all particles.}
\label{fig:cascade}
\end{figure}

\begin{figure}[h!]
\centering
\begin{subfigure}[t]{0.48\textwidth}
\includegraphics[width=\textwidth]{smc_context_benchmark/fig_pareto_time_kl.pdf}
\caption{Pareto front: time vs posterior quality.}
\label{fig:pareto}
\end{subfigure}
\hfill
\begin{subfigure}[t]{0.48\textwidth}
\includegraphics[width=\textwidth]{smc_context_benchmark/fig_n_iters.pdf}
\caption{Number of SMC iterations to convergence.}
\label{fig:n_iters}
\end{subfigure}
\caption{(a) Pareto front of total time vs final KL divergence. LSA is the reference; Hilbert+Swap achieves comparable quality in less time. (b) Methods with approximate assignments sometimes require more iterations because they accept suboptimal particles.}
\end{figure}

\paragraph{Cascade dynamics.} The progressive cascade ($\sigma_t \to \text{approx} \to \text{LSA}$) works as follows: at early iterations (large $\varepsilon$), the $\sigma_t$ acceptance rate is high (30--50\%) and the approximate methods recover many of the remaining rejects. As $\varepsilon$ decreases, the $\sigma_t$ acceptance rate drops to near zero, and the approximate step also fails to recover particles, so everything falls through to the LSA fallback. This means the cascade is most beneficial at early iterations but converges toward pure LSA cost at later iterations.


%% ====================================================================
\section{Sequential Assignment: Hamming Distance Analysis}
\label{sec:hamming}
%% ====================================================================

The key question for the sequential setting is: \emph{how different is the optimal assignment $\sigma_{t+1}^*$ from the previous optimal $\sigma_t^*$?} If consecutive optima are close in Hamming distance, a local refinement (swap) starting from $\sigma_t^*$ should suffice.

\subsection{Setup}

We run full LSA-based permABC-SMC and extract the optimal assignment $\sigma_t^*$ at each iteration. For consecutive pairs $(\sigma_t^*, \sigma_{t+1}^*)$, we compute:
\begin{enumerate}
  \item \textbf{Normalized Hamming distance}: $d_H(\sigma_t^*, \sigma_{t+1}^*)/K = \frac{1}{K}\#\{k : \sigma_t^*(k) \neq \sigma_{t+1}^*(k)\}$
  \item \textbf{Swap recovery}: starting from $\sigma_t^*$, apply pairwise swap refinement ($O(K^2)$, 3 sweeps) on the cost matrix at time $t{+}1$, then check if the result matches $\sigma_{t+1}^*$.
  \item \textbf{Na\"ive reuse}: distance ratio when simply reusing $\sigma_t^*$ at time $t{+}1$ without any refinement.
\end{enumerate}

Parameters: Gaussian model, $N=500$, $n_{\text{obs}}=10$, $K \in \{5, 10, 20, 50\}$, 5 repetitions, seed 42.

\subsection{Results}

\begin{table}[h!]
\centering
\caption{Hamming distance between consecutive optima and swap warm-start recovery (mean over all iterations and repeats).}
\label{tab:hamming}
\begin{tabular}{l cccc}
\toprule
$K$ & $\bar{d}_H / K$ & $\Pr(\sigma_t^* = \sigma_{t+1}^*)$ & Swap exact match & Swap ratio \\
\midrule
5   & 0.29 & 63\% & \textbf{86\%} & 0.989 \\
10  & 0.30 & 65\% & \textbf{86\%} & 0.989 \\
20  & 0.29 & 67\% & \textbf{84\%} & 0.995 \\
50  & 0.44 & 54\% & 59\% & 1.004 \\
\bottomrule
\end{tabular}
\end{table}

\paragraph{Interpretation.}
\begin{itemize}
  \item The normalized Hamming distance between consecutive optimal assignments remains below 0.45 for all $K$ (Table~\ref{tab:hamming}, column 2). Remarkably, 54--67\% of particles retain exactly the same optimal permutation between consecutive iterations.

  \item Starting from $\sigma_t^*$, the $O(K^2)$ swap refinement recovers the \emph{exact} LSA optimum $\sigma_{t+1}^*$ for \textbf{84--86\% of particles} when $K \le 20$. Even when the exact permutation is not recovered, the distance ratio remains $\le 1.005$, meaning the quality loss is negligible.

  \item Na\"ive reuse of $\sigma_t^*$ without refinement yields distance ratios of 2.2--3.6$\times$ (Table~\ref{tab:hamming_reuse}), confirming that while consecutive optima are \emph{close}, they are not close enough to skip refinement entirely.

  \item For $K = 50$, the swap exact match drops to 59\%, justifying a progressive cascade (swap $\to$ LSA fallback) for large $K$.
\end{itemize}

\begin{table}[h!]
\centering
\caption{Comparison: swap warm-start vs na\"ive reuse of $\sigma_t^*$ (distance ratio $d/d_{\text{LSA}}$).}
\label{tab:hamming_reuse}
\begin{tabular}{l cc}
\toprule
$K$ & Swap from $\sigma_t^*$ (ratio) & Na\"ive reuse of $\sigma_t^*$ (ratio) \\
\midrule
5   & 0.989 & 2.50 \\
10  & 0.989 & 2.23 \\
20  & 0.995 & 2.28 \\
50  & 1.004 & 3.62 \\
\bottomrule
\end{tabular}
\end{table}


\subsection{Why Swap Outperforms Hilbert in the Sequential Setting}

The micro benchmark (Section~\ref{sec:micro}) shows that Hilbert alone produces assignments with Hamming distance $\sim 0.95$ to the LSA optimum (it finds a \emph{completely different} permutation). Its distance ratio is 1.01--1.04, which is good but not optimal.

By contrast, swap warm-start from $\sigma_t^*$ exploits the continuity of the SMC trajectory: since the proposal kernel makes small perturbations to $\theta$, the optimal assignment changes by only a few transpositions. Swap is perfectly suited to find these local corrections, achieving both lower Hamming distance ($< 0.05$ in the positions it actually changes) and better distance ratio (0.989--1.004).

\textbf{The swap is therefore the best sequential assignment strategy for $K \le 20$}: it is faster than LSA ($O(K^2)$ vs $O(K^3)$), achieves near-optimal quality via warm-start, and does not require any auxiliary data structure (unlike Hilbert's curve or Sinkhorn's kernel matrix). For $K \ge 50$, the swap exact match rate degrades, and a progressive cascade (smart swap) or Hilbert+Swap is recommended.


\begin{figure}[h!]
\centering
\includegraphics[width=\textwidth]{hamming_analysis/fig_appendix_hamming_swap.pdf}
\caption{Hamming distance analysis of consecutive optimal assignments in permABC-SMC.
  \textbf{(a)} The normalized Hamming distance $d_H(\sigma_t^*, \sigma_{t+1}^*)/K$ remains below 0.45, and 54--67\% of particles keep the same optimal permutation.
  \textbf{(b)} Swap from $\sigma_t^*$ recovers the exact LSA optimum for 84--86\% of particles ($K \le 20$), with distance ratio $\le 1.005$.
  \textbf{(c)} Hamming distance vs SMC progress: it tends to decrease as the posterior concentrates.
  \textbf{(d)} Swap warm-start (solid) achieves ratio $\approx 1.0$ throughout, while na\"ive reuse of $\sigma_t^*$ (dashed) yields ratio 2--5$\times$.}
\label{fig:hamming_appendix}
\end{figure}


%% ====================================================================
\section{Updated Conclusions}
\label{sec:conclusions}
%% ====================================================================

Combining the micro benchmark (Section~\ref{sec:micro}), SMC-context results (Section~\ref{sec:smc_context}), and Hamming analysis (Section~\ref{sec:hamming}), we reach the following conclusions:

\subsection{Assignment Method Hierarchy}

\begin{enumerate}
  \item \textbf{LSA (scipy)} remains the gold standard: exact, fast enough for $K \le 50$ ($< 25$ms per iteration for $N=500$), and no quality compromise. The custom C solver with warm-start provides no measurable speedup.

  \item \textbf{Swap warm-start from $\sigma_t^*$ is the best sequential compromise} for $K \le 20$:
  \begin{itemize}
    \item $5{-}15\times$ faster than LSA
    \item Recovers the exact LSA optimum 84--86\% of the time
    \item Distance ratio $\le 1.005$ even when not exact
    \item Works because $d_H(\sigma_t^*, \sigma_{t+1}^*)/K \approx 0.29{-}0.30$
  \end{itemize}

  \item \textbf{Hilbert+Swap} is the best non-sequential alternative: $1.5{-}2.5\times$ faster than LSA, $<0.3\%$ excess cost. Useful as a fallback for $K \ge 50$ or as the initial assignment (no $\sigma_{t-1}$ available).

  \item \textbf{Smart Swap} ($\sigma_t \to \text{Swap} \to \text{LSA}$) combines the best of both: swap handles 84--86\% of particles cheaply, LSA handles the rest exactly. The cascade is most effective at early SMC iterations; at convergence, most particles fall through to LSA.

  \item \textbf{Sinkhorn} is not competitive: slower than LSA even with Numba optimization, and quality degrades at $K \ge 20$ (5--8\% excess). Interesting as a theoretical link to entropic OT, but not a practical assignment strategy for permABC.

  \item \textbf{Hilbert alone} produces permutations that are far from the LSA optimum in Hamming distance ($\sim 0.95$), even though the distance ratio is only 1--5\% worse. This means Hilbert assignments are qualitatively different from LSA, which matters when the permutation itself carries meaning (e.g., for label identification in compartmental models).
\end{enumerate}

\subsection{Practical Recommendations for the Paper}

\begin{center}
\begin{tabular}{lll}
\toprule
$K$ & Recommended method & Justification \\
\midrule
$K \le 10$  & Swap from $\sigma_t^*$ & Near-optimal, $10\times$ faster \\
$10 < K \le 50$ & Smart Swap or Hilbert+Swap & Swap alone: 59\% exact at $K=50$ \\
$K > 50$ & LSA (scipy) or Smart H+S & Swap less reliable, LSA still fast \\
\midrule
First iteration & LSA or Hilbert+Swap & No $\sigma_{t-1}$ available \\
\bottomrule
\end{tabular}
\end{center}

\end{document}
